% GENERATED FILE. Edit canonical lessons, then run npm run generate:books.
% canonical-source-hash: fnv1a64:a07225770806f4de
% canonical-lessons: TE-C312-katha, TE-C312-ata, TE-C312-sinima, TE-C312-siksakudu, TE-C312-injaniru

\chapter{Story, Game, Film, Coach, Engineer}
\label{ch:te-story-game-film-coach-engineer}
\hlchaptermodality{312}

\begin{chapteropening}
\textbf{By the end of this chapter:} \emph{I can say a story, a game, a film, a coach and an engineer.}

Complete the last lesson of chapter 312: I can say a story, a game, a film, a coach and an engineer.
\end{chapteropening}

\section[\texorpdfstring{\te{కథ} (katha)}{కథ (katha)}]{\te{కథ} (katha) — a story}

\label{lesson:TE-C312-katha}



\begin{quote}
Before the new one: say the Telugu for experience, then the Telugu for an account.
\end{quote}

\subsection*{You'll want to know: \te{కథ}}
\textbf{\te{కథ}} — \emph{katha} — \textquotedblleft{}a story\textquotedblright{}.

\begin{grammarlens}[title={What you've built}]
One more word for this chapter.
\end{grammarlens}

\subsection*{Your turn}
\begin{itemize}
\raggedright
  \item \emph{Say it:} \emph{katha}
  \item \emph{Say it:} \emph{katha}, once more
  \item \emph{Say it:} say \emph{katha}
  \item \emph{Recall:} say the Telugu for experience, then the Telugu for an account, then say \emph{katha} again
\end{itemize}

\begin{tcolorbox}[breakable,colback=teal!4,colframe=teal!35!black,title={Before you move on}]
What does \te{కథ} mean? (\textquotedblleft{}A story\textquotedblright{}.) Say it once more.
\end{tcolorbox}
\section[\texorpdfstring{\te{ఆట} (āṭa)}{ఆట (āṭa)}]{\te{ఆట} (āṭa) — a game, a sport}

\label{lesson:TE-C312-ata}



\begin{quote}
Before the new one: say the Telugu for an account, then the Telugu for a story.
\end{quote}

\subsection*{You'll want to know: \te{ఆట}}
\textbf{\te{ఆట}} — \emph{āṭa} — \textquotedblleft{}a game, a sport\textquotedblright{}.

\begin{grammarlens}[title={What you've built}]
One more word for this chapter.
\end{grammarlens}

\subsection*{Your turn}
\begin{itemize}
\raggedright
  \item \emph{Say it:} \emph{āṭa}
  \item \emph{Say it:} \emph{āṭa}, once more
  \item \emph{Say it:} say \emph{āṭa}
  \item \emph{Recall:} say the Telugu for an account, then the Telugu for a story, then say \emph{āṭa} again
\end{itemize}

\begin{tcolorbox}[breakable,colback=teal!4,colframe=teal!35!black,title={Before you move on}]
What does \te{ఆట} mean? (\textquotedblleft{}A game, a sport\textquotedblright{}.) Say it once more.
\end{tcolorbox}
\section[\texorpdfstring{\te{సినిమా} (sinimā)}{సినిమా (sinimā)}]{\te{సినిమా} (sinimā) — a film, the cinema}

\label{lesson:TE-C312-sinima}



\begin{quote}
Before the new one: say the Telugu for a story, then the Telugu for a game.
\end{quote}

\subsection*{You'll want to know: \te{సినిమా}}
\textbf{\te{సినిమా}} — \emph{sinimā} — \textquotedblleft{}a film, the cinema\textquotedblright{}.

\begin{grammarlens}[title={What you've built}]
One more word for this chapter.
\end{grammarlens}

\subsection*{Your turn}
\begin{itemize}
\raggedright
  \item \emph{Say it:} \emph{sinimā}
  \item \emph{Say it:} \emph{sinimā}, once more
  \item \emph{Say it:} say \emph{sinimā}
  \item \emph{Recall:} say the Telugu for a story, then the Telugu for a game, then say \emph{sinimā} again
\end{itemize}

\begin{tcolorbox}[breakable,colback=teal!4,colframe=teal!35!black,title={Before you move on}]
What does \te{సినిమా} mean? (\textquotedblleft{}A film, the cinema\textquotedblright{}.) Say it once more.
\end{tcolorbox}
\section[\texorpdfstring{\te{శిక్షకుడు} (śikṣakuḍu)}{శిక్షకుడు (śikṣakuḍu)}]{\te{శిక్షకుడు} (śikṣakuḍu) — a coach, a trainer}

\label{lesson:TE-C312-siksakudu}



\begin{quote}
Before the new one: say the Telugu for a game, then the Telugu for a film.
\end{quote}

\subsection*{You'll want to know: \te{శిక్షకుడు}}
\textbf{\te{శిక్షకుడు}} — \emph{śikṣakuḍu} — \textquotedblleft{}a coach, a trainer\textquotedblright{}.

\begin{grammarlens}[title={What you've built}]
One more word for this chapter.
\end{grammarlens}

\subsection*{Your turn}
\begin{itemize}
\raggedright
  \item \emph{Say it:} \emph{śikṣakuḍu}
  \item \emph{Say it:} \emph{śikṣakuḍu}, once more
  \item \emph{Say it:} say \emph{śikṣakuḍu}
  \item \emph{Recall:} say the Telugu for a game, then the Telugu for a film, then say \emph{śikṣakuḍu} again
\end{itemize}

\begin{tcolorbox}[breakable,colback=teal!4,colframe=teal!35!black,title={Before you move on}]
What does \te{శిక్షకుడు} mean? (\textquotedblleft{}A coach, a trainer\textquotedblright{}.) Say it once more.
\end{tcolorbox}
\section[\texorpdfstring{\te{ఇంజనీరు} (iñjanīru)}{ఇంజనీరు (iñjanīru)}]{\te{ఇంజనీరు} (iñjanīru) — an engineer}

\label{lesson:TE-C312-injaniru}



\begin{quote}
Before the new one: say the Telugu for a film, then the Telugu for a coach.
\end{quote}

\subsection*{You'll want to know: \te{ఇంజనీరు}}
\textbf{\te{ఇంజనీరు}} — \emph{iñjanīru} — \textquotedblleft{}an engineer\textquotedblright{}.

\begin{grammarlens}[title={What you've built}]
One more word for this chapter.
\end{grammarlens}

\subsection*{Your turn}
\begin{itemize}
\raggedright
  \item \emph{Say it:} \emph{iñjanīru}
  \item \emph{Say it:} \emph{iñjanīru}, once more
  \item \emph{Say it:} say \emph{iñjanīru}
  \item \emph{Recall:} say the Telugu for a film, then the Telugu for a coach, then say \emph{iñjanīru} again
\end{itemize}

\begin{tcolorbox}[breakable,colback=teal!4,colframe=teal!35!black,title={Before you move on}]
What does \te{ఇంజనీరు} mean? (\textquotedblleft{}An engineer\textquotedblright{}.) Say it once more.
\end{tcolorbox}
